% Lösungen Einheit 3 – Nr. 40–49
\einheitenkopf[le3][3 Exponentialfunktion]{Einheit 3 von 5 · Exponentialfunktion aufstellen und auswerten}

\erg{40}{a) Start 2020, gesucht 2025 \quad b) Start 2018, gesucht 2030 \quad c) Start 2024, gesucht 2031 \quad d) Start 2022, gesucht 2027}
\erg{41}{a) $N_0 = 100$, $q = 2$, $N(t) = 100 \cdot 2^t$ \quad b) $50$, $3$, $N(t) = 50 \cdot 3^t$ \quad c) $400$, $1{,}1$, $N(t) = 400 \cdot 1{,}1^t$ \quad d) $2500$, $1{,}02$, $N(t) = 2500 \cdot 1{,}02^t$ \quad e) $90$, $0{,}8$, $N(t) = 90 \cdot 0{,}8^t$ \quad f) $12$, $0{,}985$, $N(t) = 12 \cdot 0{,}985^t$}
\erg{42}{a) $y = 50 \cdot 2^x$ \quad b) $y = 600 \cdot 0{,}95^x$ \quad c) $y = 3000 \cdot 1{,}012^x$}
\erg{43}{a) 300: Zellen zu Beginn; 2: Verdopplung je Stunde; $t$: Zeit in Stunden \quad b) 18\,000: Neupreis in €; 0,85: jedes Jahr bleiben 85\,\%, also 15\,\% Wertverlust; $t$: Alter in Jahren \quad c) 100\,\% des alten Werts plus 4\,\% Zuwachs sind 104\,\% = 1,04 \quad d) 50: Masse bei der Geburt in kg; 1,03: jede Woche nimmt die Masse um 3\,\% zu; $x$: Zeit in Wochen nach der Geburt}
\erg{44}{a) 40 \quad b) 90 \quad c) 121 \quad d) 125 \quad e) $\approx 539{,}73$ \quad f) $t = 8$: $4000 \cdot 1{,}015^8 \approx 4506$ Einwohner \quad g) $24\,000 \cdot 0{,}88^5 \approx 12\,665{,}57\,€$}
\erg{45}{a) $N(4) = 1518{,}75 > 1500$ – stimmt \quad b) B: $5000 \cdot 1{,}02^{10} \approx 6095$; B sagt etwa 95 Einwohner mehr voraus als A \quad c) $A(t) = 30 \cdot 1{,}25^t$; $A(21) \approx 3252{,}6\,$m² $> 3000\,$m² – die Behauptung stimmt}
\erg{46}{a) $t = 4$, $N = 160$ \quad b) $t = 4$, $N = 243$ \quad c) $t = 4$, $N \approx 1220{,}70$ \quad d) $t = 6$, $N \approx 94{,}12$ \quad e) 2028: $\approx 971\,485{,}52\,€$ (noch darunter); 2029: $420\,000 \cdot 1{,}15^7 \approx 1\,117\,208{,}35\,€$ – im Jahr 2029}
\erg{47}{a) Jahreszahl statt Zahl der Schritte eingesetzt; von 2020 bis 2026 sind es 6 Schritte: $600 \cdot 1{,}03^6 \approx 716$ Tiere \quad b) Faktor mit der Schrittzahl multipliziert statt potenziert: $1200\,€ \cdot 1{,}04^5 \approx 1459{,}98\,€$ \quad c) Prozentsatz statt Faktor als Basis: $N(t) = 90 \cdot 1{,}04^t$}
\erg{48}{a) Jedes Jahr wird der neue Wert wieder mit 1,1 multipliziert: $500 \cdot 1{,}1 \cdot 1{,}1 \cdot 1{,}1 = 500 \cdot 1{,}1^3$; „mal 3“ würde den Wert verdreifachen. \quad b) Beim zweiten und dritten Mal werden die 10\,\% vom schon erhöhten Wert genommen: $1{,}1^3 = 1{,}331$, also 33,1\,\% statt 30\,\%.}
\erg{49}{a) $A(t) = 40\,000 \cdot 1{,}015^t$, \ $B(t) = 46\,000 \cdot 0{,}99^t$ \quad b) $A(5) \approx 43\,091$, \ $B(5) \approx 43\,746$ \quad c) nach 6 Jahren ($A(6) \approx 43\,738$, $B(6) \approx 43\,308$)}
